\documentclass[11pt,a4paper]{article}

\usepackage[utf8]{inputenc}
\usepackage[T1]{fontenc}
\usepackage{lmodern}
\usepackage{geometry}
\usepackage{graphicx}
\usepackage{booktabs}
\usepackage{array}
\usepackage{amsmath}
\usepackage{amssymb}
\usepackage{hyperref}
\usepackage{natbib}
\usepackage{xcolor}
\usepackage{sectsty}
\usepackage{caption}
\usepackage{enumitem}

\geometry{margin=1in}
\sectionfont{\sffamily\bfseries\large\MakeUppercase}
\subsectionfont{\sffamily\bfseries\normalsize}
\captionsetup{font=small,labelfont=bf}
\setlist{itemsep=0.2em, topsep=0.3em}

\definecolor{abstractblue}{RGB}{218,226,255}

\title{\sffamily\bfseries\LARGE Evaluating Open-Vocabulary Vision-Language Models for 2D Robotic Grasp Rectangle Detection}
\author{
Zhenkun Pang\\
\small MSc Robotics \& Artificial Intelligence\\
\small Supervisor: Dr Jan Paul Siebert
}
\date{}

\begin{document}

\maketitle

\section*{Abstract}
\noindent\colorbox{abstractblue}{%
\parbox{\dimexpr\linewidth-2\fboxsep\relax}{%
This project investigates whether pre-trained open-vocabulary vision-language models can support the visual perception stage of 2D robotic grasp detection. The proposed pipeline uses the Cornell Grasping Dataset, applies open-vocabulary object grounding to localise target objects from text prompts, and converts the grounded regions into oriented grasp rectangles using classical computer vision geometry. The study compares a traditional OpenCV-based baseline with a VLM-assisted grasp detection pipeline using localisation IoU, grasp rectangle IoU, angular error, Cornell-style grasp success criteria, and qualitative failure analysis. Optional few-shot adaptation will be explored only if the main pipeline is stable.
}}

\vspace{1em}
\noindent\textbf{Keywords:} robotic grasp detection, vision-language models, open-vocabulary grounding, Cornell Grasping Dataset, computer vision

\section{Introduction}

Robotic grasping depends on reliable visual perception. A robot must identify a target object, estimate its shape and orientation, and predict a feasible grasp configuration. In many 2D grasp detection benchmarks, this configuration is represented as an oriented grasp rectangle.

Recent vision-language models have shown strong open-vocabulary recognition and grounding capabilities. These models can localise image regions from natural-language prompts, making them attractive for robotic perception tasks where objects may not belong to a fixed predefined taxonomy. However, it remains unclear whether their object grounding outputs are directly useful for robotic grasp detection.

This project studies a focused question: can an open-vocabulary vision-language model act as a perception front-end for 2D robotic grasp rectangle detection when combined with classical geometric post-processing? The project does not aim to build a full robot manipulation system. Instead, it evaluates the visual perception component using an offline benchmark dataset.

\section{Research Questions}

The central research question is:

\begin{quote}
To what extent can pre-trained open-vocabulary vision-language models, when combined with classical geometric grasp estimation, support 2D robotic grasp rectangle detection?
\end{quote}

The project will address four sub-questions:

\begin{enumerate}
    \item Can pre-trained VLMs or open-vocabulary detection models localise objects in the Cornell Grasping Dataset using text prompts without task-specific end-to-end training?
    \item Can VLM-localised regions be converted into reliable oriented grasp rectangles using OpenCV-based methods such as contour extraction, PCA, and rotated bounding boxes?
    \item How does a VLM-assisted pipeline compare with a traditional computer vision baseline in grasp accuracy, angle error, and robustness?
    \item If time and resources allow, can lightweight few-shot adaptation improve performance on difficult object categories?
\end{enumerate}

\section{Methods and Materials}

\subsection{Dataset}

The project will use the Cornell Grasping Dataset, a common benchmark for 2D robotic grasp detection. It contains RGB-D images of everyday objects and manually annotated grasp rectangles. The dataset is suitable for an MSc project because it is public, widely used, and allows quantitative evaluation without requiring physical robot hardware.

\subsection{Grasp Representation}

A 2D grasp rectangle can be represented as:

\begin{equation}
    g = (x, y, w, h, \theta),
\end{equation}

where $(x, y)$ is the grasp centre, $w$ and $h$ describe the gripper-related rectangle dimensions, and $\theta$ is the grasp orientation.

\subsection{Traditional Computer Vision Baseline}

The baseline will use classical image processing and geometric estimation. Candidate methods include thresholding, contour extraction, morphological filtering, PCA-based orientation estimation, and minimum-area rotated bounding boxes. This baseline provides a reproducible reference system and a fallback if model deployment becomes unstable.

\subsection{VLM-Assisted Pipeline}

The VLM-assisted pipeline will use a pre-trained open-vocabulary grounding model such as Grounding DINO, OWL-ViT, or Florence-2. The model will receive image inputs and text prompts describing the target object. Its predicted object region will then be passed to the same geometric grasp estimation backend used by the baseline.

\subsection{Evaluation}

The planned evaluation metrics include:

\begin{itemize}
    \item localisation IoU for target object regions;
    \item grasp rectangle IoU against Cornell annotations;
    \item angular error between predicted and annotated grasp orientations;
    \item Cornell-style grasp success criteria;
    \item qualitative analysis of failure cases.
\end{itemize}

\section{Proposed Pipeline}

\begin{figure}[h]
    \centering
    \fbox{\parbox{0.88\linewidth}{
    \centering
    \vspace{1em}
    Cornell RGB/RGB-D image + text prompt\\
    $\downarrow$\\
    Open-vocabulary object grounding\\
    $\downarrow$\\
    Region cropping or segmentation cue\\
    $\downarrow$\\
    OpenCV geometric grasp estimation\\
    $\downarrow$\\
    Oriented 2D grasp rectangle\\
    $\downarrow$\\
    Quantitative and qualitative evaluation
    \vspace{1em}
    }}
    \caption{Overview of the planned VLM-assisted grasp detection pipeline.}
    \label{fig:pipeline}
\end{figure}

\section{Project Plan}

\begin{table}[h]
\centering
\caption{Initial project timeline.}
\label{tab:timeline}
\begin{tabular}{p{0.18\linewidth}p{0.28\linewidth}p{0.42\linewidth}}
\toprule
Time & Stage & Main milestone \\
\midrule
Weeks 1--2 & Literature and dataset preparation & Read grasp detection and VLM grounding papers; parse Cornell images and annotations; visualise ground-truth grasps. \\
Weeks 3--4 & CV baseline & Implement OpenCV segmentation, contour analysis, PCA or rotated bounding boxes, and basic metrics. \\
Weeks 5--6 & VLM-assisted pipeline & Deploy a grounding model and connect object localisation to geometric grasp estimation. \\
Weeks 7--8 & Testing and optional adaptation & Run experiments across object categories; attempt lightweight adaptation only if the main system is stable. \\
Weeks 9--10 & Robustness and failure analysis & Compare baseline and VLM-assisted results; analyse localisation, geometry, and angle failures. \\
Weeks 11--12 & Dissertation writing & Finalise figures, tables, discussion, references, and submission formatting. \\
\bottomrule
\end{tabular}
\end{table}

\section{Expected Contributions}

The expected contributions are:

\begin{itemize}
    \item a reproducible Cornell dataset processing and evaluation pipeline;
    \item a traditional computer vision baseline for 2D grasp rectangle prediction;
    \item a VLM-assisted grasp detection pipeline combining open-vocabulary grounding and geometric grasp estimation;
    \item a quantitative comparison using standard grasp detection metrics;
    \item a failure-mode analysis of VLM-based grounding in robotic grasp perception.
\end{itemize}

\section{Risk Management}

The highest-risk component is optional few-shot adaptation or parameter-efficient fine-tuning, because it may require additional compute, model-specific engineering, and environment setup. The safe core deliverable is therefore an offline grasp detection evaluation pipeline on the Cornell dataset, including dataset parsing, visualisation, a CV baseline, VLM-assisted inference, metrics, and failure analysis.

If fine-tuning is not feasible, the project will remain focused on zero-shot VLM-assisted grasp detection. If VLM deployment becomes unstable, the project will prioritise the traditional CV baseline and evaluation framework.

\section{Acknowledgments}

This project proposal was prepared for the MSc Robotics \& Artificial Intelligence dissertation planning stage.

\bibliographystyle{plainnat}
\bibliography{references}

\end{document}
